\begin{textsample}{POS Dim 3 – se – Score 29.00 – se/se\_inf\_1.txt}  \label{ex:f3_pos_161}
2.
A ETAPA EDUCAÇÃO \textbf{INFANTIL} > Certa vez, quando tinha seis anos, vi num \textbf{livro} sobre a Floresta Virgem [... ] uma imponente gravura.
Representava ela uma jiboia que engolia uma fera : dizia o \textbf{livro} : ‘ as jiboias engolem, sem mastigar a presa inteira.
Em seguida não podem mover -se e dormem os seis \textbf{meses} da digestão ’.
Refleti então sobre as aventuras na selva, e fiz, com meu \textbf{lápis} de cor, o meu primeiro desenho[... ]. mostrei às pessoas grandes e perguntei se meu desenho lhes fazia medo.
Responderam -me : ‘ por que um chapéu faria medo? ’ Meu desenho não representava um chapéu.
Representava uma jiboia digerindo um elefante.
Desenhei então, o interior da jiboia, a fim de que as pessoas grandes pudessem compreender.
Elas têm sempre a necessidade de explicações [... ].
As pessoas grandes aconselharam -me a deixar de lado o desenho de jiboias abertas ou fechadas, e dedicar -me de preferência, à geografia, à história, ao cálculo, à gramática.
Foi assim que abandonei, aos seis anos, uma esplêndida carreira de pintor.
Eu fora desencorajado pelo insucesso do desenho número 1 e do desenho número 2.
As pessoas grandes não compreendem nada sozinhas, e é cansativo para as \textbf{crianças}, estar a toda hora explicando. ( SAINT-EXUPÉRY, 1994, p.7-8 ) O texto preliminar apresentado encaminha -se na perspectiva de fomentar o desencadeamento das contribuições para a elaboração do documento Curricular do Estado de Sergipe, no que concerne à etapa da Educação \textbf{Infantil}.
Este processo abre importante espaço para participação dos profissionais da educação, pois, entendemos o currículo como uma construção coletiva, arraigada ao contexto que dialoga com os sujeitos de cada comunidade.
Portanto, as especificidades do nosso Estado devem ser alçadas como ponto de partida para um projeto educativo, que promova às novas gerações de sergipanos a conquista de conhecimentos socialmente úteis para os contextos de vida nos quais estão inseridas.
Na elaboração do currículo dessa etapa, foi priorizada a \textbf{criança} como sujeito do mundo.
Assim, optamos por reafirmar a identidade das \textbf{crianças} como sujeitos históricos e de direitos que passam por processos de aprendizagem de acordo com as especificidades do seu desenvolvimento, à luz dos preceitos das Diretrizes Curriculares Nacionais para a Educação \textbf{Infantil} ( Parecer CNE/CEB no 20/09 ; Res.
CNE/CEB no 05/09 ) e da BNCC, que trazem os Campos de Experiências e os Direitos de Aprendizagem e Desenvolvimento.
Para melhor contribuir com o conhecimento de todos os docentes, o texto introdutório do documento curricular ficou organizado com a seguinte estrutura : - Contexto Histórico-político da Educação \textbf{Infantil} no Brasil ; - Fundamentos Filosófico-Pedagógicos da Educação \textbf{Infantil} ; - Elementos para a Organização Curricular na Educação \textbf{Infantil}.
Nesta perspectiva, \textbf{convidamos} todos os profissionais da educação e comunidade em geral para que realizem a leitura desse documento preliminar, discutam e apresentem \textbf{sugestões} que venham a contribuir com a construção do Currículo Sergipano. 2.1 Contexto Histórico-Político da Educação \textbf{Infantil} no Brasil 2.1.1 Aspectos da trajetória histórica da Educação \textbf{Infantil} A trajetória da Educação \textbf{Infantil} brasileira tem sua história marcada pelas múltiplas concepções de “ \textbf{criança} ”, de “ \textbf{infância} ” e de “ desenvolvimento \textbf{infantil} ” estando alinhada à história dos homens. > [... ] a história das \textbf{instituições} pré-escolares não é uma sucessão de fatos que se somam, mas a interação de tempos, influências e temas, em que o período de elaboração da proposta educacional assistencialista se integra aos outros tempos da história dos homens ( KUHLMANN, Jr., 2000, p.81 ).
Foram os interesses e as influências políticas e sociais de cada período histórico que delinearam as necessidades e as formas de atendimento às \textbf{crianças}, a exemplo do ocorrido na segunda metade do século XIX quando o atendimento à \textbf{infância} adquiriu um caráter assistencialista, em virtude do conjunto de transformações políticas, econômicas e sociais vivenciadas na Europa e, consequentemente, no Brasil, a exemplo da Revolução Industrial. > [... ] as \textbf{creches} e \textbf{pré-escolas} surgiram a partir de mudanças econômicas, políticas e sociais que ocorreram na sociedade : pela incorporação das mulheres à força de trabalho assalariado, na organização das famílias, num novo papel da mulher, numa nova relação entre os sexos, para citar apenas as mais evidentes.
Mas, também, por razões que se identificam com um conjunto de ideias novas sobre a \textbf{infância}, sobre o papel da \textbf{criança} na sociedade e de como torná -la, através da educação, um indivíduo produtivo e ajustado às exigências desse conjunto social ( BUJES, 2001, p.15 ).
Além da Revolução Industrial, a Proclamação da República brasileira ( 1889 ) trouxe efeitos econômicos, culturais e sociais que modificaram a estrutura de vida da população, intensificando a necessidade da família em possuir apoio para os \textbf{cuidados} com suas proles, uma vez que a presença feminina no mercado industrial tomou proporção considerável. > Enquanto para as famílias mais abastadas pagavam uma babá, as pobres se viam na contingência de deixar os filhos sozinhos ou colocá -los numa \textbf{instituição} que deles \textbf{cuidasse}.
Para os filhos das mulheres trabalhadoras, a \textbf{creche} tinha que ser de tempo integral ; para os filhos de operárias de baixa renda, tinha que ser gratuita ou cobrar muito pouco ; ou para \textbf{cuidar} da \textbf{criança} enquanto a mãe estava trabalhando fora de casa, tinha que zelar pela saúde, ensinar \textbf{hábitos} de \textbf{higiene} e alimentar a \textbf{criança}.
A educação permanecia assunto de família.
Essa origem determinou a associação \textbf{creche}, \textbf{criança} pobre e o caráter assistencial da \textbf{creche} ( DIDONET, 2001, p. 13 ).
Assim, \textbf{creches}, jardins de \textbf{infância} e \textbf{parques} \textbf{infantis} foram criados, e, junto a eles, eclodiu a necessidade de regulamentação do atendimento.
Nessa conjuntura, surgiram iniciativas de atendimento à \textbf{infância} com o propósito hospitaleiro e higienista de subsídio aos extratos sociais desfavorecidos, visando ainda afastar as \textbf{crianças} do trabalho servil, além de servirem como guardiãs de \textbf{crianças} órfãs.
As \textbf{creches} ainda foram vistas como substituição ou oposição à \textbf{roda} dos expostos “ para que as mães não abandonassem suas crianças”[1 ] ( KUHLMANN JR,1999, p. 82 ).
Naquela conjuntura histórica, acreditava -se que as \textbf{crianças} das camadas populares eram consideradas “ [... ] carentes, deficientes e inferiores na medida em que não correspondem ao padrão estabelecido ; faltariam a essas \textbf{crianças} privadas culturalmente, determinados atributos ou conteúdos que deveriam ser nelas incutidos ” ( KRAMER, 1995, p. 24 ).
Para compensar tal carência, a proposta educacional para essas \textbf{crianças} diferenciava -se da proposta oferecida nas \textbf{instituições} onde estavam as \textbf{crianças} de famílias mais abastadas.
Nessa perspectiva, a \textbf{pré-escola} funcionaria, segundo a Kramer ( 1995 ), como mola propulsora da mudança social, uma vez que possibilitaria a democratização das oportunidades educacionais.
A proposta pedagógica das unidades de \textbf{pré-escola} seguia a influência dos modelos pedagógicos europeus e americanos, traduzida pelo movimento da Escola Nova, o qual teve à frente Fernando de Azevedo, Anísio Teixeira e Lourenço Filho.
Entre as concepções, era defendida a “ escola pública organizada em sistemas de ensino, desde o jardim de \textbf{infância} até a universidade e trabalhavam com a ideia de educação popular, na primeira fase da escolarização ” ( MIGUEL, 2004, p.31 ).
A diferença da proposta ainda era mais perceptível nos estabelecimentos privados cujo cunho pedagógico estava voltado à socialização, criatividade e à preparação para o ensino regular ( KRAMER, 1995 ).
Essa diferença no atendimento com base na origem social foi dirimida com a implementação de políticas governamentais, a exemplo da Lei no 5.692/71, a qual assegura em seu Art. 19, § 2o, que “ os sistemas velarão para que as \textbf{crianças} de idade inferior a 7 anos recebam educação em escolas maternais, jardins-de-infância ou \textbf{instituições} equivalentes ”.
Já a Constituição Federal de 1988 garantiu no art. 208, inciso IV que “ o dever do Estado para com a educação será efetivado mediante a garantia de oferta de \textbf{creches} e \textbf{pré-escolas} às \textbf{crianças} de zero a seis anos de idade ” ( BRASIL, 1988 ).
As \textbf{creches}, por conseguinte, passam para a responsabilidade da educação com ênfase nas funções pedagógicas fundadas nos aspectos cognitivos da \textbf{criança} e na garantia de direitos.
A Constituição Federal foi “ um marco decisivo na afirmação dos direitos da \textbf{criança} no Brasil ” ( LEITE FILHO, 2001, p. 31 ).
Outro marco foi o Estatuto da Criança e do Adolescente ( 1990 ) que em seu art. 3o determina que : > A \textbf{criança} e o adolescente gozam de todos os direitos fundamentais inerentes à pessoa humana, sem prejuízo da proteção integral de que trata esta Lei, assegurando -se-lhes, por lei ou por outros meios, todas as oportunidades e facilidades, a fim de lhes facultar o desenvolvimento físico, mental, moral, espiritual e social, em condições de liberdade e de dignidade ( BRASIL, 1990 ).
Nas palavras de Ferreira, o ECA garantiu o direito das \textbf{crianças}, a exemplo do “ direito ao \textbf{afeto}, direito de \textbf{brincar}, direito de querer, direito de não querer, direito de conhecer, direito de sonhar.
Isso quer dizer que são atores do próprio desenvolvimento ” ( 2000, p. 184 ).
Seguindo na trajetória dos marcos legais, destaca -se a Lei no 9.394/96, Lei das Diretrizes e Bases da Educação Nacional que, ao tratar da composição dos níveis escolares, inseriu a educação \textbf{infantil} como primeira etapa da Educação Básica, com matrícula obrigatória a partir dos 4 anos de idade, com redação dada pela Lei no 12.796/2013 ; e considerando o desenvolvimento integral da \textbf{criança} em “ seus aspectos físico, psicológico, intelectual e social, complementando a ação da família e da comunidade ” ( BRASIL, 1996, art. 29 ).
Complementando a legislação, o Ministério da Educação instituiu as Diretrizes Curriculares Nacionais para a Educação \textbf{Infantil} ( 2009 ), sendo este um instrumento orientador da organização das atividades cotidianas das \textbf{instituições} de Educação \textbf{Infantil}.
Destaca -se ainda o Referencial Curricular da Educação \textbf{Infantil} que atende as expectativas de aprendizagem e desenvolvimento da \textbf{criança} do século XXI na perspectiva da educação integral.
Por fim, destacamos a Base Nacional Comum da Educação \textbf{Infantil} e do Ensino Fundamental, que traz um novo olhar para a Educação \textbf{Infantil}, contemplando os seis direitos de aprendizagens, os quais dialogam com os campos de experiências e com as competências gerais.
A \textbf{criança} foi pensada como um sujeito que aprende de forma integral, nas mais diversas situações da vida a partir das interações e \textbf{brincadeiras}. 2.1.2 Aspectos da Educação \textbf{Infantil} em Sergipe Em 1931, o Interventor Federal Augusto Maynard encarregou o secretário da Diretoria Geral da Instrução Pública do Estado, professor José Augusto da Rocha Lima, para visitar as \textbf{instituições} escolares de São Paulo a fim de estudar os novos métodos pedagógicos com objetivo de implementá -los em Sergipe.
O seu relatório apontou as influências da escola ativa e os ideais da Escola Nova com incentivo da espontaneidade criadora da \textbf{criança} a partir das teorias de Decroly, Dewey, Rousseau, Pestalozzi e Froebel, para o favorecimento da prática pedagógica ( ANDRADE, 1931 ).
Nesse mesmo ano, Dr.
Helvécio de Andrade estava à frente da Diretoria Geral da Instrução Pública, momento em que escreveu um Relatório Anual apresentando ao Interventor Federal Augusto Maynard suas ações na educação pública estadual.
Em seus escritos ressaltou que : > Certo, o início do ensino do pré-escolar jardim escolar, no próximo ano, com inauguração do jardim de \textbf{infância}, a cuja organização de V.
Exa. dedicado todos os seus esforços, trará para o ensino público sergipano uma fase de aperfeiçoamento digna de todos os aplausos.
Na primeira idade, que vai até os cinco anos, os movimentos \textbf{infantis} são incoordenados, despersivos.
Disciplinar esses movimentos, dar -lhes uma significação, transformá -los em \textbf{hábitos} salutares, de conduta, de \textbf{higiene}, de observação, é lançar as bases de uma educação completa ( ANDRADE, 1931, p. 2).[2 ] No ano seguinte à escrita do Relatório Anual, precisamente em 17 de março de 1932, foi inaugurada a primeira escola de Educação \textbf{Infantil}, denominada Casa da Criança de Sergipe, cujo nome, posteriormente, foi alterado para Jardim de Infância Augusto Maynard, em homenagem ao Interventor Federal que atendeu os anseios das professoras normalistas Helena Abud e Miriam Santos Melo ( FERNANDEZ, s/d ).
O Jardim de Infância teve como primeira Diretora, a \textbf{Professora} laranjeirense Penélope Magalhães dos Santos, a qual buscou inspirações pedagógicas nas correntes que estavam em prática no Rio de Janeiro e em São Paulo ( \textbf{NASCIMENTO}, s.d, p. 5 ).
Ressalta -se que o jardim de Infância Augusto Maynard marcou a trajetória educacional de Sergipe, funcionando de 1932 a 2002.
Essa \textbf{instituição} foi construída para \textbf{acolher} as \textbf{crianças} de 04 a 06 anos de idade, com preocupação voltada à saúde, à alimentação, à \textbf{higiene} e à sobrevivência.
No interior do estabelecimento, existia uma Inspetoria de \textbf{Higiene} \textbf{Infantil} cuja premissa era a proteção à \textbf{infância}, \textbf{cuidados} de saúde e \textbf{higiene}, além do preparo físico como requisito para o ingresso escolar ( ANDRADE, 1931 ).
Na década de 40, a realidade educacional do Estado era alarmante : > Sergipe entrou nesse período com 72,6 por cento da sua população analfabeta.
Existiam apenas 536 escolas públicas primárias ( escolas isoladas com apenas uma única sala ), sendo 358 estaduais e 178 municipais, com uma matrícula total de 21.810 alunos.
Havia ainda cerca de 120 escolas particulares e apenas 2 \textbf{instituições} de educação \textbf{infantil} – o Jardim “ Augusto Maynard ” e a escola da Legião da Boa \textbf{Vontade} ( LBA ) ( LIMA, 2006, p. 158).[3 ] Do exposto, verifica -se que muitas \textbf{crianças} e jovens ainda estavam fora da escola, quadro esse que sofreu \textbf{pequena} alteração com a inauguração, em 1944, do primeiro Jardim de Infância Municipal de Aracaju, denominado de Centro Municipal de Assistência à Criança, “ constituído de um Jardim de Infância e de um Posto Puericultura, uma \textbf{Biblioteca} \textbf{infantil} e Grupo Escolar ” ( Decreto no 75, de 21 de março de 1945, art. 1o ).
O Centro visava ao preparo de \textbf{crianças} para a “ escola primária, de 4 a 6 anos, em 2 anos, divididos em períodos de 8 \textbf{meses}, incluídas as férias, por exercícios apropriados que visem a desenvolver o espírito de observação e formar \textbf{hábitos} mentais sociais e higiênicos ” ( Decreto no 75, de 21 de março de 1945, art. 3o ).
Esta \textbf{instituição}, situada no Bairro Siqueira Campos, ficou conhecida como o Jardim Operário porque fora a primeira \textbf{instituição} de atendimento pré-escolar pública municipal situada num bairro popular, que atendia \textbf{crianças} das famílias das camadas trabalhadoras ( DECRETO LEI no 75, 1945, apud LIMA, 2006, p. 158 ).
Embora o Decreto Lei no 75/1945 instituísse o atendimento aos filhos das camadas trabalhadoras do bairro Siqueira Campos, acabou por beneficiar as classes privilegiadas : > A matrícula no Jardim de Infância do Centro Municipal de Assistência à Criança obedecia a alguns pré-requisitos que sustentavam a ideia de uma \textbf{instituição} para alguns privilegiados, os que tinham boas condições e status social baseadas nos padrões éticos, morais e de saúde perfeita.
Não era qualquer \textbf{criança} que poderia matricular -se no Jardim de Infância (... ).
As exigências para o ingresso na \textbf{instituição} eram claras no sentido de que não poderiam se matricular \textbf{crianças} com princípio de leitura e cálculo.
Visto que, assegura -se a ideia de que era o Jardim que iria cumprir o papel de ensinar as primeiras letras ( SANTOS ; CARDOSO ; MELO, 2016, p. 6 ).
Como se observa do exposto, não podiam se matricular no Jardim de Infância \textbf{crianças} que já “ tivessem princípios de leitura e de cálculo, que fosse portador de defeito físico, moléstias contagiosas e predisposições mórbidas de caráter degenerativo ” ( DECRETO no 98, 1932 apud LEAL, s.d., p. 4 ).
As construções de Jardins de Infância em Sergipe representaram um relevante avanço no que se refere à instrução e aos \textbf{cuidados} com a \textbf{higiene} e saúde ; representaram ainda o impulso das reformas educacionais, a exemplo do que ocorria em vários Estados brasileiros.
Na década de 1960, surgiram unidades de ensino vinculadas ao Sistema “ S ” “ criadas para ofertar a Educação \textbf{Infantil} aos filhos dos trabalhadores, a exemplo do Jardim de Infância Pequeno Polegar e Jardim de Infância João Bolinha, ambos pertencentes ao Serviço Social da Indústria – SESI ” ( PEE, 2015, p. 6 ).
Os jardins de \textbf{infância} ou escolas \textbf{infantis} originaram -se no âmbito público governamental, com algumas iniciativas particulares também.
Sua expansão foi “ lenta e gradual até os anos 1970, apesar de um início de \textbf{crescimento} nos anos 1950 com a criação das classes de pré-primário, anexas aos estabelecimentos de ensino fundamental ” ( SANTOS ; CARDOSO ; MELO, 2016, p. 8 ).
Nos anos de 1960 e 1970, verificou -se uma expansão das \textbf{creches} e jardins de \textbf{infância} de natureza privada e assistencial, devido ao aumento da taxa do trabalho feminino.
Os anos de 1980 foram marcados pelas dimensões legais no ordenamento jurídico a partir da Constituição de 1988, que buscou inspirações na Declaração Universal dos Direitos da Criança e do Adolescente ( 1959 ).
Pela Constituição Federal é dever : > da família, da sociedade e do Estado assegurar à \textbf{criança}, ao adolescente e ao jovem, com absoluta prioridade, o direito à vida, à saúde, à alimentação, à educação, ao lazer, à profissionalização, à cultura, à dignidade, ao respeito, à liberdade e à convivência familiar e comunitária, além de colocá -los a salvo de toda forma de negligência, discriminação, exploração, violência, crueldade e opressão ( Art. 227 ).
O inciso IV, do art. 228, da CF, define a “ educação \textbf{infantil}, em \textbf{creche} e \textbf{pré-escola}, às \textbf{crianças} até 5 ( cinco ) anos de idade ”.
Em Sergipe, seguindo o que preceitua o inciso V, art. 11, da Lei no 9.394/96, o atendimento às \textbf{crianças} de \textbf{creches} e \textbf{pré-escolas} foi \textbf{transferido} para os municípios, sendo o atendimento da rede estadual finalizado no ano de 2010 ( PEE, 2015, 8 ).
Observa -se ainda o \textbf{crescimento} do atendimento das redes privadas entre os anos de 2007 e 2010, conforme constam nos gráficos abaixo : Imagem 01 – quadro de estabelecimentos que ofertam \textbf{creche} em Sergipe por dependência administrativa.
Fonte : INEP/MEC/Sinopses Estatísticas apud PEE, 2015, p. 8.
Imagem 02 – quadro de estabelecimentos que ofertam \textbf{pré-escola} em Sergipe por dependência administrativa.
Fonte : INEP/MEC/Sinopses Estatísticas apud PEE, 2015, p. 8.
Dessa análise estatística, inserida no diagnóstico realizado no Plano Estadual de Educação-PEE, em 2015, depreende -se que : > No período analisado ( 2007 a 2013 ), observamos que o número de \textbf{creches} na rede municipal cresceu em 30,89 \%, enquanto a rede privada apresentou \textbf{crescimento} muito maior, na ordem de 165,96 \%.
Quando observamos o número de estabelecimentos que ofertam a \textbf{pré-escola}, percebemos que a rede privada cresceu 43,29 \% e a rede municipal teve um decréscimo de 7,5 \%.
Os dados indicam que a rede privada vem investindo nessa etapa de ensino de modo mais acelerado que o poder público, embora a rede pública possua o maior número de estabelecimentos ( p. 8 ).
Pelo diagnóstico do PEE, no período de 2009 a 2013, “ a Educação \textbf{Infantil} apresentou um \textbf{crescimento} de 2,60 \% nas matrículas : foram 5.308 matrículas a mais na \textbf{creche} e 3.326 na \textbf{pré-escola} ” ( 2015, p. 9 ).
No que se refere ao atendimento das \textbf{creches}, 81 \% das matrículas concentram -se nas redes municipais de ensino, ficando 19 \% para rede privada ( PARENTE ; PARENTE, 2010 ).
Segundo dados da Plataforma QEdu, no ano de 2017, o atendimento de Educação \textbf{Infantil} ocorreu em 1.496 escolas públicas e privadas, tanto da zona rural quanto da zona urbana.
Esses estabelecimentos atenderam a 22.347 \textbf{crianças} de \textbf{creche} e 58.995 \textbf{crianças} na \textbf{pré-escola}, perfazendo um total de 81.342 matrículas.
Esse atendimento da rede pública ainda será ampliado com a finalização das construções de \textbf{instituições} de ensino com incentivos advindos do Programa Proinfância, do Governo Federal.
As perspectivas de ampliação do atendimento, previstas no Plano Estadual de Educação e nos Planos Municipais de Educação, devem ser acompanhadas de ações em prol da organização de um currículo sergipano que expresse, de forma qualificada, as concepções pedagógicas voltadas ao desenvolvimento e aprendizagem das \textbf{crianças} de 0 a 5 anos de idade. [ 1 ] : A \textbf{roda} dos expostos, como assistência caritativa, era, pois, missionária.
A primeira preocupação do sistema para com a \textbf{criança} nela deixada era de providenciar o batismo, salvando a alma da \textbf{criança}, a menos que trouxesse consigo um bilhete – o que era muito comum – que informava à rodeira de que o \textbf{bebê} já estava batizado.
No caso de dúvida dos responsáveis pela \textbf{instituição}, a \textbf{criança} era novamente batizada.
Mas o fenômeno de abandonar os filhos é tão antigo como a história da colonização brasileira, só que antes da \textbf{roda}, as \textbf{crianças} eram abandonadas e supostamente assistidas pelas municipalidades, ou pela compaixão de quem as encontrava ( AQUINO, 2001, p. 31 ). [ 2 ] : A ortografia foi atualizada. [ 3 ] : Quadro Demonstrativo dos Grupos Escolares, Escolas Reunidas, Escolas Isoladas, Estaduais, Municipais e Particulares e Nomes dos Respectivos Professores.
Governo de Sergipe.
Imprensa Oficial.
Aracaju, Abril, 1941.

% matched lemmas: acolher, afeto, bebê, biblioteca, brincadeira, brincar, convidar, creche, crescimento, criança, cuidado, cuidar, higiene, hábito, infantil, infância, instituição, livro, lápis, mês, nascimento, parque, pequeno, professora, pré-escola, roda, sugestão, transferir, vontade
\end{textsample}
